\documentclass[11pt]{article}
\ifdefined\XeTeXrevision\else\ifdefined\pdfoutput\pdfoutput=1\fi\fi

\usepackage[margin=1in]{geometry}
\usepackage[T1]{fontenc}
\usepackage{amsmath,amssymb,amsthm}
\usepackage{booktabs}
\usepackage{tikz}
\usepackage{xcolor}
\definecolor{linkblue}{RGB}{25,60,120}
\usepackage[colorlinks=true,linkcolor=linkblue,citecolor=linkblue,urlcolor=linkblue]{hyperref}
\hypersetup{
  pdftitle={Achievement numbers of four small polyominoes},
  pdfauthor={Jan Philipp Harries},
  pdfsubject={Polyomino achievement games},
  pdfkeywords={polyomino achievement game, Harary's generalized tic-tac-toe, Maker-Breaker game, board number, move number, computer-assisted proof, certificate, OEIS A380597, OEIS A380598}
}

\newtheorem{theorem}{Theorem}
\newtheorem{lemma}[theorem]{Lemma}
\newtheorem{proposition}[theorem]{Proposition}
\newtheorem{corollary}[theorem]{Corollary}
\theoremstyle{definition}
\newtheorem{definition}[theorem]{Definition}
\theoremstyle{remark}
\newtheorem{remark}[theorem]{Remark}

\newcommand{\bweak}{b^{\mathrm w}}
\newcommand{\mweak}{m^{\mathrm w}}
\newcommand{\bstr}{b^{\mathrm s}}
\newcommand{\mstr}{m^{\mathrm s}}
\newcommand{\file}[1]{\texttt{#1}}
\newcommand{\Pl}{\mathcal{P}}

% one polyomino drawn from a list of cells (x,y), y upwards
\newcommand{\poly}[1]{%
  \begin{tikzpicture}[scale=0.32,baseline=0]
    \foreach \x/\y in {#1} {\filldraw[fill=gray!35,draw=black] (\x,\y) rectangle ++(1,1);}
  \end{tikzpicture}}

\title{Achievement numbers of four small polyominoes}
\author{Jan Philipp Harries\thanks{ellamind GmbH, Bremen, Germany. \textit{E-mail address:} \href{mailto:jan@ellamind.com}{jan@ellamind.com}}}
\date{October 8, 2026}

\begin{document}
\maketitle

\begin{abstract}
In Harary's achievement game for a polyomino $P$ two players alternately claim cells of an
$n\times n$ board, and the first player tries to own a congruent copy of $P$. Gardner's tables give,
for each polyomino that the first player can achieve, a board number $b$ (the side of the smallest
square board on which the first player wins) and a move number $m$ (the number of the first
player's moves needed there); OEIS A380597 and A380598 reproduce the values of Gardner's 2001 book.
For the I-tetromino and the L-, Y- and N-pentominoes we determine $(b,m)$ in the Maker--Breaker
(weak) game as $(6,6)$, $(6,7)$, $(6,6)$ and $(5,6)$, and in the strong game, in which a player who
completes a copy first wins, as $(7,7)$, $(7,9)$, $(6,6)$ and $(5,6)$. Hence the tabulated pairs
$(7,9)$ for the Y-pentomino and the board number $6$ for the N-pentomino are wrong in both games;
for the I-tetromino, tabulated $(7,8)$, and the L-pentomino, tabulated $(7,10)$, both values are
wrong in the weak game, and in the strong game the board number $7$ is correct and the move number
is wrong. Every statement is proved by an explicit strategy that short programs check
exhaustively; for the strong game on $6\times 6$, and for the move numbers on $7\times 7$ of the
I-tetromino and the L-pentomino, these are strategies of the second player that block, complete
copies of their own, and end lines in pairing strategies. The strategies and programs are in the ancillary files.
\end{abstract}

\section{Introduction}

Harary's generalization of tic-tac-toe, popularized by Gardner in 1979~\cite{Gar79}, fixes a
polyomino (an ``animal'') $P$ and a square board. Two players alternately claim free cells, and a
player who owns all cells of a copy of $P$---rotations and reflections allowed---wins. By strategy
stealing the second player cannot win with correct play, so the game is a first-player win or a
draw. Gardner describes this \emph{strong} game: the first player to complete the animal wins, and
for a polyomino that cannot be achieved the second player ``can always force a draw''
\cite[p.~497]{Gar01}. Later work, for example~\cite{HSP07,Sie04}, studies the \emph{weak} or Maker--Breaker version, in
which the second player only blocks.

For every polyomino that the first player can achieve, Gardner tabulates the \emph{board number}
$b$, the side of the smallest square board on which the first player can force a win, and the
\emph{move number} $m$, the number of the first player's moves needed to win on that board. The
tables in Gardner's 1992 and 2001 books~\cite{Gar92,Gar01} (in~\cite{Gar01}: chapter~37,
Figures~37.2 and~37.3) give no proofs, and Gardner invites the reader to prove the
values~\cite[p.~497]{Gar01}. They differ in several entries from the 1979 column as quoted
in~\cite{Wiki}, whose table had already been questioned in~\cite{Tho96}. \emph{Winning Ways}~\cite{BCG82}
apparently tabulates a different pair, the least number of moves on an unbounded board and the least
board on which that number suffices~\cite{Tho96}; we do not compare with it.
The OEIS entries A380597 and A380598~\cite{OEIS597,OEIS598}, by P.~von Br\"omssen (2025), list $b$
and $m$ for the free polyominoes in the order of A246521~\cite{OEIS521}; for every polyomino that
the first player can achieve they agree with Gardner's 2001 tables. Neither entry says which game
is meant.

This note settles the board and move numbers of four of these polyominoes in both games
(Section~\ref{sec:results}). The other entries of size at most five
agree with the tables in the weak game (Remark~\ref{rem:small}). The proofs are finite: each
statement is reduced to an explicit strategy for one of the players, written as a certificate file
and checked exhaustively by programs that share no code with the search that found it
(Section~\ref{sec:method}).

\begin{figure}[ht]
\centering
\begin{tabular}{ccccc}
\poly{0/0,1/0,2/0,3/0} &
\poly{0/0,1/0,2/0,3/0,0/1} &
\poly{0/0,1/0,2/0,3/0,1/1} &
\poly{0/1,1/1,1/0,2/0,3/0} &
\poly{0/0,1/0,2/0,3/0,3/1,4/1}\\[2pt]
I-tetromino ($I_4$) & L-pentomino ($L_5$) & Y-pentomino ($Y_5$) & N-pentomino ($N_5$) & Snaky\\
$n=9$ & $n=13$ & $n=14$ & $n=15$ & $n=45$
\end{tabular}
\caption{The polyominoes of this note and their indices $n$ in A380597/A380598 (the polyomino with
code A246521$(n+1)$).}
\label{fig:shapes}
\end{figure}

\section{Games, board numbers and move numbers}\label{sec:defs}

Let $[n]=\{0,\dots,n-1\}$ and $B_n=[n]^2$. A \emph{placement} of a polyomino $P\subset\mathbb Z^2$
on $B_n$ is a set $g(P)\subseteq B_n$, where $g$ is one of the eight symmetries of the square
lattice followed by a translation; $\Pl_n(P)$ is the set of placements. Two players, $X$ (first)
and $O$, alternately claim free cells of $B_n$, starting with $X$.

\begin{definition}
In the \emph{weak game} $X$ (Maker) wins as soon as he owns a placement; $O$ (Breaker) wins if this
never happens. In the \emph{strong game} a player who owns a placement before the other does wins;
if the board fills first, the game is a draw. For $d\ge 1$, \emph{$X$ wins within $d$ moves} if $X$
has a strategy such that against every play of $O$ he owns a placement after at most $d$ of his own
moves (in the strong game: and $O$ owns no placement before that moment). The \emph{board number}
$\bweak(P)$ is the least $n$ such that $X$ wins the weak game on $B_n$ within some $d$, and the
\emph{move number} $\mweak(P)$ is the least $d$ such that $X$ wins the weak game on $B_{\bweak(P)}$ within
$d$ moves. $\bstr(P)$ and $\mstr(P)$ are defined in the same way with the strong game.
\end{definition}

This is the convention of Gardner's tables: for the I-tromino the board number is $4$ and the move
number $3$~\cite[p.~494]{Gar01}. The following two facts are standard; we include proofs because
the lower bounds below rest on them.

\begin{lemma}\label{lem:weakstrong}
Fix $P$, $n$ and $d$.
\begin{enumerate}
\item[(a)] If $O$ has a strategy in the weak game on $B_n$ against which $X$ owns no placement after
any of his first $d$ moves, then $X$ does not win the strong game on $B_n$ within $d$ moves.
\item[(b)] If $X$ wins the strong game on $B_n$ within $d$ moves, he wins the weak game on $B_n$
within $d$ moves.
\end{enumerate}
\end{lemma}

\begin{proof}
(a) In the strong game $O$ follows the blocking strategy. As long as the strong game goes on, the
positions are positions of the weak game under that strategy, so $X$ owns no placement after any of
his first $d$ moves; if the strong game stops earlier, it stops with a draw or a win for $O$.
(b) $X$ uses the same strategy. Until $X$ completes a placement the two games pass through the same
positions, and the strong strategy completes one within $d$ moves.
\end{proof}

\begin{lemma}\label{lem:monotone}
If $k\le n$ and $X$ wins the weak game on $B_k$ within $d$ moves, then $X$ wins the weak game on
$B_n$ within $d$ moves. Consequently a blocking strategy for $O$ against $d$ moves (or against any
number of moves) on $B_n$ yields one on every $B_k$ with $k\le n$, since finite games of perfect information are
determined.
\end{lemma}

\begin{proof}
Let $\sigma$ be the winning strategy on $B_k\subseteq B_n$. $X$ plays $\sigma$ in an imagined game on
$B_k$ in which $X$ owns the same cells as in reality and $O$ owns a superset of $O$'s real cells in
$B_k$. A real move of $O$ on a cell of $B_k$ that is free in the imagined game is also $O$'s
imagined move; any other move of $O$ is replaced in the imagined game by a move on an arbitrary
free cell of $B_k$. Such a cell exists as long as $X$ has not won, since $\sigma$ wins before the
imagined board is full. Every move of $\sigma$ is free in the imagined game and therefore in
reality, and every placement on $B_k$ is a placement on $B_n$, so $X$ completes one within $d$
moves.
\end{proof}

Lemma~\ref{lem:weakstrong} gives $\bweak\le\bstr$, and $\mweak\le\mstr$ when $\bweak=\bstr$. Lemmas
\ref{lem:weakstrong}(a) and~\ref{lem:monotone} turn one blocking strategy on $B_{b-1}$ into the
lower bound $b$ on the board number in both games. We use no analogue of
Lemma~\ref{lem:monotone} for the strong game: a strategy of $O$ in the strong game on $B_n$ that is
not a blocking strategy is used only on $B_n$ itself.

\section{Results}\label{sec:results}

\emph{Proved} below means that the statement follows from Lemmas~\ref{lem:weakstrong}
and~\ref{lem:monotone}, Propositions~\ref{prop:breaker} and~\ref{prop:ostrat} and from certificates
listed in Table~\ref{tab:evidence}, each accepted by every checker named there
(Section~\ref{sec:method}). \emph{Computed} means a verdict of a search
program for which no certificate exists.

\begin{theorem}[weak game]\label{thm:weak}
$(\bweak,\mweak)$ equals $(6,6)$ for $I_4$, $(6,7)$ for $L_5$, $(6,6)$ for $Y_5$ and $(5,6)$ for $N_5$.
\end{theorem}

\begin{theorem}[strong game]\label{thm:strong}
$(\bstr,\mstr)$ equals $(7,7)$ for $I_4$, $(7,9)$ for $L_5$, $(6,6)$ for $Y_5$ and $(5,6)$ for $N_5$.
In particular, for $I_4$ and for $L_5$ the first player does not win the strong game on any board
$B_n$ with $n\le 6$; with correct play of the second player such a game ends in a draw or in a win
for the second player.
\end{theorem}

\begin{proof}[Proof of Theorems~\ref{thm:weak} and~\ref{thm:strong}]
Take $I_4$ in the weak game. The Breaker certificate for the full game on $B_5$ and
Lemma~\ref{lem:monotone} give $\bweak\ge 6$; the Maker certificate on $B_6$ with budget $6$ gives
$\bweak\le 6$ and $\mweak\le 6$; the Breaker certificate against $5$ moves on $B_6$ gives $\mweak\ge 6$. The
other shapes are the same with the files of Table~\ref{tab:evidence}. In the strong game the two
Breaker certificates of a shape give, by Lemmas~\ref{lem:weakstrong}(a) and~\ref{lem:monotone}, $\bstr\ge\bweak$ and, for $Y_5$
and $N_5$, that $X$ does not win on $B_{\bweak}$ within $\mweak-1$ moves; the strong strategies give the
matching upper bounds. For $I_4$ and $L_5$: on $B_n$ with $n\le 5$, $X$ does not win the strong game
by the Breaker certificate on $B_5$ and Lemmas~\ref{lem:weakstrong}(a) and~\ref{lem:monotone}; on
$B_6$ he does not win by the $O$-strategy certificate with budget $18$ and
Proposition~\ref{prop:ostrat}, since on $B_6$ the first player makes at most $18$ moves; on $B_7$ he
wins within $7$ ($I_4$) and $9$ ($L_5$) moves by the strong strategies, and not within $6$ and $8$
moves by the $O$-strategy certificates with these budgets. Hence $\bstr=7$ and $\mstr=7$,
respectively $9$.
\end{proof}

The smallest certificate is a single pairing: on $B_4$, cells $(x,y)$, the six pairs
$\{(x,0),(x,1)\}$ for $x\in[4]$ and $\{(0,2),(1,2)\}$, $\{(2,2),(3,2)\}$ meet every placement of $N_5$
in a full pair, so $O$, answering each move of $X$ on a paired cell with its partner, prevents $N_5$
on $B_4$ (and, by Lemma~\ref{lem:monotone}, on smaller boards).

\begin{table}[ht]
\centering\small
\begin{tabular}{@{}lllrll@{}}
\toprule
shape & statement & file in \file{anc/certs/} & size & checked by\\
\midrule
$I_4$ & $O$ blocks on $B_5$ & \file{n9\_b5\_full} & 10{,}130 & A, B\\
      & $X$ wins weak game on $B_6$ within 6 & \file{n9\_b6\_moves} & 2{,}615 & A, R\\
      & $O$ blocks 5 moves on $B_6$ & \file{n9\_b6\_lt6} & 2{,}227 & A, B\\
      & $X$ wins strong game on $B_7$ within 7 & \file{strong\_n9\_b7\_d7} & 333{,}113 & S, T\\
      & $X$ does not win strong game on $B_6$ & \file{i4\_b6\_full} & 4{,}154{,}214 & C, D, E\\
      & \quad nor on $B_7$ within 6 & \file{i4\_b7\_K6} & 21{,}474 & C, D, E\\
$L_5$ & $O$ blocks on $B_5$ & \file{n13\_b5\_full} & 321 & A, B\\
      & $X$ wins weak game on $B_6$ within 7 & \file{n13\_b6\_moves} & 11{,}075 & A, R\\
      & $O$ blocks 6 moves on $B_6$ & \file{n13\_b6\_lt7} & 981 & A, B\\
      & $X$ wins strong game on $B_7$ within 9 & \file{strong\_n13\_b7\_d9} & 948{,}483 & S, T\\
      & $X$ does not win strong game on $B_6$ & \file{l5\_b6\_full} & 471{,}778 & C, D, E\\
      & \quad nor on $B_7$ within 8 & \file{l5\_b7\_K8} & 884{,}056 & C, D, E\\
$Y_5$ & $O$ blocks on $B_5$ & \file{n14\_b5\_full} & 293 & A, B\\
      & $X$ wins weak game on $B_6$ within 6 & \file{n14\_b6\_moves} & 3{,}913 & A, R\\
      & $O$ blocks 5 moves on $B_6$ & \file{n14\_b6\_lt6} & 213 & A, B\\
      & $X$ wins strong game on $B_6$ within 6 & \file{strong\_n14\_b6\_d6} & 1{,}023{,}521 & S, T\\
$N_5$ & $O$ blocks on $B_4$ & \file{n15\_b4\_full} & 1 & A, B\\
      & $X$ wins weak game on $B_5$ within 6 & \file{n15\_b5\_moves} & 3{,}882 & A, R\\
      & $O$ blocks 5 moves on $B_5$ & \file{n15\_b5\_lt6} & 84 & A, B\\
      & $X$ wins strong game on $B_5$ within 6 & \file{strong\_n15\_b5\_d6} & 253{,}534 & S, T\\
\bottomrule
\end{tabular}
\caption{Certificates behind Theorems~\ref{thm:weak} and~\ref{thm:strong}. ``Blocks'' means a
Breaker strategy in the weak game, against the full game or against the stated number of Maker
moves. ``Does not win'' is shown by a strategy of $O$ in the strong game (Section~\ref{sec:method});
the budget of the $B_6$ certificates is $18$, the whole game. Size: certificate nodes; for strong
strategies of $X$, positions with $X$ to move; for strategies of $O$, table lines, i.e.\ positions
with $O$ to move plus pairing leaves, each up to the eight symmetries of the board. The four
strategies of $O$ are stored gzipped with sizes $16.2$, $0.10$, $1.8$ and $3.7$\,MB. Checkers:
A~\file{check\_cert.py}, B~\file{fresh\_checker/check\_breaker2.py}, R~\file{play\_cert.py}
(exhaustive concrete replay), S~\file{check\_strong.py}, T~\file{fresh\_checker/check\_strong2.py},
C~\file{strong/check6.py}, D~\file{strong/check6b.cpp}, E~\file{strong/chk.py}.}
\label{tab:evidence}
\end{table}

\begin{corollary}\label{cor:oeis}
In both games the values $(b,m)=(7,9)$ tabulated for $Y_5$ are wrong (both are $6$), and the board
number $6$ tabulated for $N_5$ is wrong (it is $5$); the tabulated move number $6$ of $N_5$ is
correct. For $I_4$, tabulated $(7,8)$, and $L_5$, tabulated $(7,10)$, both values are wrong in the
weak game. In the strong game the board number $7$ is correct for both, and the move numbers are
wrong: they are $7$ for $I_4$ and $9$ for $L_5$.
\end{corollary}

In particular the example line of A380597, that the I-tetromino is a draw on square boards of side
less than $7$, is true in the strong game (Theorem~\ref{thm:strong}) and false in the weak game,
where Maker wins on $6\times 6$.

\begin{table}[ht]
\centering\small
\begin{tabular}{@{}llllll@{}}
\toprule
shape & $n$ & Gardner 1979$^{\ast}$ & Gardner 1992/2001, OEIS & weak game & strong game\\
\midrule
$I_4$ & 9  & $6$ / $8$ & $7$ / $8$  & $6$ / $6$ & $7$ / $7$\\
$L_5$ & 13 & $6$ / $9$ & $7$ / $10$ & $6$ / $7$ & $7$ / $9$\\
$Y_5$ & 14 & $6$ / --  & $7$ / $9$  & $6$ / $6$ & $6$ / $6$\\
$N_5$ & 15 & --        & $6$ / $6$  & $5$ / $6$ & $5$ / $6$\\
\bottomrule
\end{tabular}
\caption{Board number / move number. $^{\ast}$Not seen by the author; values as quoted
in~\cite{Wiki}, which annotates only some 1979 entries (``--'': no 1979 value quoted there). The weak and strong columns are Theorems~\ref{thm:weak} and~\ref{thm:strong}.}
\label{tab:compare}
\end{table}

The 1979 board number $6$ for $I_4$, $L_5$ and $Y_5$ agrees with the weak game; in the strong game
it holds for $Y_5$ and does not hold for $I_4$ and $L_5$, whose strong board number is $7$, as in the
1992/2001 tables. The 1979 pairs $(6,8)$ for $I_4$ and $(6,9)$ for $L_5$ match neither game: on
$6\times 6$ the first player wins the weak game within $6$ and $7$ moves and does not win the strong
game. We
know of no source that gives board number $5$ for $N_5$. For $I_4$ a weak-game
win on $7\times 7$ within $7$ moves was given earlier by Earnest and Mathias~\cite{EM22} (cited
in~\cite{Wiki}); Theorem~\ref{thm:weak} improves it to $6\times 6$ within $6$ moves.

\begin{remark}[computed consistency checks]\label{rem:computed}
The following solver verdicts are not used in any proof; they agree with the certified values.
(i) For $I_4$ and $L_5$ on $B_7$, the first and the second solver of Section~\ref{sec:method} both
find a win for $X$ in the strong game within $7$ and $9$ moves and none within $6$ and $8$ moves. On
$B_6$ the second solver finds no win for $X$ within $11$ moves for $I_4$ (151\,s) and within $13$
moves for $L_5$ (878\,s); the first solver reaches $8$ and $10$. The second solver does not decide
the full $6\times 6$ game within $1700$\,s; it is decided by the $O$-strategy certificates of Table~\ref{tab:evidence},
which the third solver \file{strong/solver\_strong.cpp} produced.
(ii) As controls, the second solver gives $(5,4)$ for the T-tetromino, $(4,4)$ for the
L-tetromino and $(3,5)$ for the S-tetromino in both games, as in the tables.
\end{remark}

\begin{remark}[checked with one checker]\label{rem:small}
With the first solver, Breaker and Maker certificates of the kind in Table~\ref{tab:evidence} were
also generated for all entries $n\le 21$ (polyominoes of size at most five), and each is accepted
by \file{check\_cert.py}; no second checker was run on them. They reproduce A380597 and A380598 in
the weak game except at $n=9,13,14,15$, and for the entries with value $0$ they show that $X$ does
not win on boards up to $8\times 8$. These certificates are not in the ancillary files;
\file{validate\_small.py} regenerates and checks them in under a minute.
\end{remark}

\begin{remark}[Snaky]\label{rem:snaky}
A380597 lists Snaky, $n=45$ (Figure~\ref{fig:shapes}), as its only unknown value. Halupczok and Schlage-Puchta proved that Breaker wins the weak game on $8\times 8$, with a strategy
found by computer~\cite[Theorem~1(3), Section~2]{HSP07}. Our solver produces Breaker certificates
for this board, one for each of the $10$ classes of Maker's first cell under the symmetries of the
board; each passes \file{check\_cert.py} (the largest has $1{,}126{,}032$ nodes). This reproduces
the result as a checked certificate. By Lemmas~\ref{lem:weakstrong} and~\ref{lem:monotone}, if Snaky
is achievable on some square board in either game, its board number is at least $9$. These
certificates ($191$\,MB) are not in the ancillary files; \file{split\_run.py} regenerates and checks
them in about $40$ minutes with three processes. Whether Snaky can be achieved on a finite board in
the strong game is open, and its board number is unknown in both games; for the weak game a
first-player strategy on $17\times 17$ has been announced in a preprint~\cite{OAI26}, which this
note does not use.
\end{remark}

\section{Method}\label{sec:method}

\paragraph{Representation and search.} Cells are numbered $y\cdot W+x$; positions are pairs of
bitmasks $(M,B)$, and placements are enumerated from the shape by the eight lattice symmetries and
all translations. The search program \file{solver.cpp} answers ``does $X$ win within $d$ moves?''
by depth-first AND/OR search with a transposition table keyed by $64$-bit hashes, canonical under
the symmetries of the board, and with pruning rules listed in the ancillary README. A hash collision
or a pruning error could make a verdict wrong; therefore nothing is called proved on the strength of
a verdict. For proved results the solver exports a certificate, and the certificate is checked by
programs that use no hashing and no pruning beyond the rules proved below.

\paragraph{Breaker certificates.} Let $d$ be Maker's remaining number of moves ($d=\infty$ for the
full game). A placement $p$ is \emph{$d$-completable} at $(M,B)$ if $p\cap B=\emptyset$ and
$|p\setminus M|\le d$; a free cell is \emph{$d$-relevant} if it lies in a $d$-completable
placement. A certificate is a directed acyclic graph whose nodes are evaluated along every play from
the empty board, with a \emph{virtual} position $(M,B,d)$ attached. The checker
\file{check\_cert.py} requires at every reached node that $M$ contains no placement, and:
\begin{itemize}
\item[\texttt{m}] (Maker to move, $d\ge 1$): a child for every $d$-relevant cell $c$, at
$(M\cup\{c\},B,d-1)$; if some free cell is not $d$-relevant, also a \emph{pass child} at
$(M,B,d-1)$;
\item[\texttt{b}] (Breaker to move): a free $d$-relevant cell $c$ and a child at $(M,B\cup\{c\},d)$;
\item[\texttt{Z}, \texttt{D}] (leaves): $d=0$, or no $d$-relevant cell exists;
\item[\texttt{P}] (leaf, Maker to move): pairwise disjoint pairs of free cells such that every
$d$-completable placement contains one of the pairs.
\end{itemize}

\begin{lemma}\label{lem:perm}
If $p$ is not $d$-completable at some position, it is not completable for the then remaining budget
at any later position, also when Maker's moves are replaced by passes.
\end{lemma}

\begin{proof}
A move of $O$ only enlarges $B$. A move of $X$ lowers $|p\setminus M|$ by at most one and $d$ by
exactly one; a pass lowers only $d$. So $|p\setminus M|>d$ persists.
\end{proof}

\begin{proposition}\label{prop:breaker}
If \file{check\_cert.py} accepts a Breaker certificate with budget $d_0$, then $O$ has a strategy in
the weak game against which $X$ owns no placement after any of his first $d_0$ moves.
\end{proposition}

\begin{proof}
$O$ walks the certificate. When $X$ plays a $d$-relevant cell, $O$ moves to its child; when $X$
plays another cell $g$, $O$ moves to the pass child and calls $g$ a \emph{ghost}: the virtual
position omits ghosts from $M$ and otherwise equals the real one, with the same budget. At
\texttt{b} nodes $O$ plays the given cell.

(1) \emph{No placement through a ghost is ever completable in reality.} When $g$ was played, every
placement $p\ni g$ with $p\cap B=\emptyset$ had $|p\setminus M_{\mathrm{virt}}|>d$. If $p$ contains no
earlier ghost, its real and virtual counts agree, so $p$ was not completable in reality, and it
stays so by Lemma~\ref{lem:perm}; placements through earlier ghosts are covered by induction. In the
virtual game a ghost is a pass, so by Lemma~\ref{lem:perm} it never becomes relevant there: $O$ is
never told to play a ghost, and the cells $O$ plays are free in reality.

(2) By (1), every placement that is completable in reality avoids the ghosts and is completable in
the virtual position. Hence $X$ owns no real placement while the virtual $M$ contains none
(checked at every node); at \texttt{Z} the budget is used up; at \texttt{D} nothing is completable;
at \texttt{P} every placement that $X$ could still complete within the budget contains a pair of
really free, non-ghost cells, and from then on $O$ answers each move of $X$ in a pair with its
partner (by Lemma~\ref{lem:perm} no other placement becomes completable later).
\end{proof}

At a \texttt{P} leaf a placement with one missing cell cannot contain a pair of free cells, so the
condition also excludes immediate threats. Pairing strategies are the classical way to show that a
polyomino cannot be achieved (see~\cite{Tho96,HSP07}); here only the placements completable within
the budget need to be covered.

\paragraph{Maker certificates.} A Maker certificate gives $X$'s cell at each node and, at $O$'s
turn, explicit children for some replies and a default child for the rest; \file{check\_cert.py}
accepts a default only for replies outside the default subtree's envelope (the cells $X$ claims there
and the cells of the placements he completes there). We do
not rely on this rule: \file{play\_cert.py} follows each Maker certificate as a strategy against
\emph{every} legal reply of $O$ at every reached position, memoized on (node, $M$, $B$), and
requires a real placement within the budget. The four Maker certificates of
Table~\ref{tab:evidence} pass this replay with $174{,}766$, $702{,}984$, $288{,}636$ and $89{,}274$
distinct states.

\paragraph{Strong-game strategies.} A strong certificate is a table from positions with $X$ to move
to $X$'s move. \file{check\_strong.py} starts from the empty board, plays the table's move, stops if
$X$ owns a placement, and otherwise fails if the budget is used up or the board is full; it then
tries every free cell as $O$'s reply, fails if $O$ owns a placement, and continues. No pruning is
used, so acceptance is a complete proof that $X$ wins within the budget. The tables were exported
with every reply of $O$ enumerated; the four tables of Table~\ref{tab:evidence} have
$253{,}534$ to $1{,}023{,}521$ positions, and every table entry is reached.

\paragraph{Strategies of $O$ in the strong game.} To show that $X$ does not win the strong game on
$B_n$ within $K$ moves we give a strategy of $O$. The game ends when $O$ completes a placement first,
so $O$ may attack as well as block; a pure blocking strategy, as used for the weak game, does not
exist on $B_6$ for $I_4$ and $L_5$, because Maker wins the weak game there. For a position $(X,O)$
let $r=K-|X|$ be the number of moves $X$ has left, and call a placement $p$ \emph{$r$-live} if
$p\cap O=\emptyset$ and $|p\setminus X|\le r$.

A certificate of kind \texttt{ostrategy} (format in the ancillary README) consists of a header
(board, shape, kind, budget $K$) and a table indexed by \emph{canonical keys}: the key of $(X,O)$ is
the least pair $(g(X),g(O))$, compared as integers via bitmasks, over the eight symmetries $g$ of
the board. An entry \texttt{O} gives $O$'s move at a position with $O$ to move; an entry \texttt{P}
gives pairs of cells at a position with $X$ to move. Cells are stored in the frame of the key; at a
real position $Q$ the checker picks a symmetry $g$ with $g(Q)$ equal to the key and maps a stored
cell $c$ to $g^{-1}(c)$. The checker walks the game from the empty board. At a position with $X$ to
move it tries \emph{every} free cell and fails if the move completes a placement of $X$; at a
position with $O$ to move it requires a table entry whose cell is free. A line of play ends, and is
accepted, exactly at these positions:
\begin{itemize}
\item[(a)] $O$'s move completes a placement of $O$;
\item[(b)] $X$ has made $K$ moves ($r=0$);
\item[(c)] the board is full;
\item[(d)] no placement is $r$-live (\emph{hopeless});
\item[(e)] $X$ is to move and the table has a \texttt{P} entry whose pairs, mapped to the real
board, consist of two distinct free cells each, are pairwise disjoint, and every $r$-live placement
contains both cells of some pair (\emph{pairing leaf}).
\end{itemize}
Tests (b)--(d) are made at positions of both kinds. No other pruning is used: $X$'s moves are never
restricted. The three checkers test these conditions in different orders; since each accepted leaf
is sound on its own, the order does not matter.

\begin{proposition}\label{prop:ostrat}
If one of the three checkers in \file{strong/} (\file{check6.py}, \file{check6b.cpp},
\file{chk.py}) accepts an \texttt{ostrategy} certificate with budget $K$ on $B_n$, then $O$ has a strategy in the
strong game on $B_n$ against which $X$ never owns a placement after one of his first $K$ moves
unless $O$ owned a placement before. In particular $X$ does not win the strong game on $B_n$ within
$K$ moves, and not at all if $K\ge \lceil n^2/2\rceil$.
\end{proposition}

\begin{proof}
\emph{Symmetry.} \file{check6.py} and \file{chk.py} verify that the set
$\Pl_n(P)$ is mapped onto itself by the eight symmetries of the board, so the conditions checked at
a position---legality of a move, completion of a placement, (b)--(e)---hold at $Q$ for a move $m$ if
and only if they hold at $h(Q)$ for $h(m)$, for every symmetry $h$, and $Q$ and $h(Q)$ have the same
key. \file{check6b.cpp} needs no such argument: it memoizes on exact positions and verifies every real position that occurs,
with the move it plays there. \file{check6.py} memoizes on keys: when it meets a position $Q'$
whose key it has verified at $Q$, then $Q'=h(Q)$ for some symmetry $h$, and the verified tree
below $Q$ maps under $h$ to a verified tree below $Q'$. \file{chk.py} verifies the canonical
positions themselves and canonicalizes every child, which is the same argument. In all three cases
$O$'s strategy is: at a real position $Q$ whose key carries an \texttt{O} entry $c$, play
$g^{-1}(c)$ for a symmetry $g$ with $g(Q)$ equal to the key; different choices of $g$ differ by a
symmetry fixing $Q$ and give moves with isomorphic verified subtrees.

\emph{Internal positions.} Along any play that stays inside the verified tree, no move of $X$
completes a placement, since every free cell was tried.

\emph{Leaves.} At (a) $O$ owns a placement and $X$ none, and the game is over. At (c) the game is
a draw. At (b) $X$ has made $K$ moves without owning a placement. At (d) and (e) let $O$ continue
as follows: if $X$ claims a cell of a pair whose partner is free, $O$ claims the partner; otherwise
$O$ claims any free cell. (At (d) there are no pairs.) A placement $p$ that $X$ owns after at most
$r$ further moves is disjoint from $O$'s cells at all later times, hence at the leaf, and needs at
most $r$ further cells, so it was $r$-live at the leaf. At (d) no such placement exists. At (e)
it contains a pair $\{u,v\}$ of cells that were free at the leaf. The statement ``every pair is
entirely free or contains a cell of $O$'' holds at the leaf and is preserved: when $X$ claims a cell
of an entirely free pair, $O$ claims its partner at once (the pairs are disjoint, so one answer
suffices), and a pair containing a cell of $O$ keeps it. Hence $X$ never owns both $u$ and $v$, and
never owns $p$. If $O$ completes a placement or the board fills before,
the game ends earlier, without a win for $X$.

Finally, $X$ makes at most $\lceil n^2/2\rceil$ moves on $B_n$, so with $K\ge\lceil n^2/2\rceil$
the budget is no restriction.
\end{proof}

The pairing rule (e) covers only placements completable within the remaining budget, as in
Proposition~\ref{prop:breaker}; it is applied only at positions where it is checked directly, so no
analogue of the pass children of Breaker certificates is needed. For the $6\times 6$ certificates
the budget is $18$, the whole game, and (b) never occurs.

The $O$ strategies of Table~\ref{tab:evidence} were found by a third solver,
\file{strong/solver\_strong.cpp}: depth-first proof-number search~\cite{Nag02} with a transposition
table keyed by the full canonical position (two $64$-bit words, so without hash collisions),
immediate wins, double threats and forced blocks for both players, the leaves (d) and (e) (pairings
found by a bounded backtracking search), and a move restriction to cells in placements still
completable by either player. That restriction is not proved for the strong game and is not used by
the checkers: the export enumerates every free cell at each position with $X$ to move and solves
every position the search had pruned. After an opening move of $X$ on the border of $B_6$, $O$'s
answer already leads to a pairing leaf; after an interior opening the strategy also uses threats of its own,
and $3{,}331{,}388$ of the $3{,}880{,}764$ \texttt{O} entries of the $I_4$ certificate complete a
placement of $O$. \file{strong/check6.py} and \file{strong/chk.py} report the counts of the leaf
kinds and of unused table entries; there are no unused entries in any of the four certificates. Expanded to exact positions, the $I_4$ certificate on $B_6$ covers
$2{,}071{,}013$ positions with $X$ to move and $15{,}939{,}736$ with $O$ to move.

\paragraph{Independent programs.} The directory \file{fresh\_checker/} holds a second checker for
Breaker certificates, \file{check\_breaker2.py}, and for strong strategies,
\file{check\_strong2.py}, written independently, from the format description alone, with their own
placement enumeration; they accept every Breaker certificate and every strategy of $X$ in
Table~\ref{tab:evidence}. A second solver,
\file{fresh\_solver/ttt.cpp}, was written independently, from the rules of the game alone, with its own pruning; on boards
up to $4\times4$ it agrees with an unpruned minimax (\file{naive.cpp}) and on $5\times 5$ with an
unpruned budgeted search (\file{naive\_b.cpp}) for budgets $3$ to $5$, in both games. It reproduces
every value of Theorem~\ref{thm:weak}, the values $(6,6)$ and $(5,6)$ of Theorem~\ref{thm:strong}
and the $7\times 7$ verdicts of Remark~\ref{rem:computed}(i). The script \file{fresh\_checker/index\_own.py} enumerates
the polyominoes of size at most five, recomputes A246521 and confirms the indices $9,13,14,15$.

The three checkers for strategies of $O$ share no code with \file{strong/solver\_strong.cpp} or with
each other. \file{strong/check6.py} (Python) memoizes on canonical keys; \file{strong/check6b.cpp}
(C++) memoizes on exact positions and has its own placement enumeration and symmetry formulas;
\file{strong/chk.py} (Python) was written from the format description alone, without reading the
solver or the other two checkers. Each takes the claimed board, shape, kind and budget on the command
line and fails if the certificate header differs; each prints the number of placements ($36$ and
$120$ for $I_4$ and $L_5$ on $B_6$, $56$ and $192$ on $B_7$, the counts $2n(n-3)$ and $8(n-3)(n-1)$).
They report the same numbers of \texttt{O} and \texttt{P} table lines. In further tests, corrupted
certificates (a dropped
\texttt{O} or \texttt{P} entry, a move onto an occupied cell or off the board, a pair dropped, a pair
on an occupied cell, overlapping pairs, the budget raised by one) were rejected by all three
checkers. On boards up to $4\times 4$, and on $5\times 5$ for budgets up to $6$, the third solver's
verdicts agree with an unpruned minimax search for $14$ polyominoes, and the certificates it exports
there pass \file{strong/check6.py} and \file{strong/check6b.cpp}. The $X$ strategies of
Table~\ref{tab:evidence}, converted to canonical keys, pass these two checkers as well. The scripts
of these tests are not part of the ancillary files.

\paragraph{What is trusted.} The proofs depend on: the checkers (a few hundred lines each,
available in two separately written versions for Breaker certificates and strategies of $X$ in the
strong game, and in three for strategies of $O$) and their placement enumeration; the proof rules of
Propositions~\ref{prop:breaker} and~\ref{prop:ostrat}; the replay \file{play\_cert.py} for the four
weak Maker wins, which for $Y_5$ and $N_5$ also follow from the strong strategies by
Lemma~\ref{lem:weakstrong}(b); the Python interpreter; and, for \file{strong/check6b.cpp}, the C++
compiler. They do not depend on any of the three solvers.

\section{Reproduction}\label{sec:repro}

The ancillary directory \file{anc/} contains the certificates of Table~\ref{tab:evidence}, the
programs named above, a manifest of SHA-256 sums and a README with every command and its expected
output. With \texttt{uv} and a C++17 compiler:

{\small
\begin{verbatim}
cd anc
shasum -a 256 -c MANIFEST
bash verify_all.sh      # 44 checks, about 4 minutes; ends "ALL 44 CHECKS PASSED"

clang++ -O2 -std=c++17 -o solver solver.cpp      # regenerate, e.g. N5 strong:
./solver --board 5 --shape "1,0;0,1;2,0;1,1;3,0" --game strong --dmin 6 --dmax 6 \
    --tt-mb 512 --cert s15.txt && uv run python check_strong.py s15.txt
bash regenerate_weak.sh   # the 12 weak certificates, checked; a few seconds

cd fresh_solver && clang++ -O2 -std=c++17 -o ttt ttt.cpp
./ttt Y5 6 strong 7      # K=5 no_win_within_K, K=6 P1_WINS
./ttt L5 7 strong 10     # K=8 no_win_within_K, K=9 P1_WINS

cd ../strong            # strategies of O; note -fno-vectorize
clang++ -O2 -fno-vectorize -std=c++17 -o solver_strong solver_strong.cpp
./solver_strong --board 6 --shape "0,0;1,0;2,0;3,0" --K 18 --tt-log2 23 \
    --max-rss-mb 3800 --cert i4_b6_full.txt          # DONE x_cannot_win ...
uv run python chk.py i4_b6_full.txt --board 6 --shape "0,0;1,0;2,0;3,0" \
    --kind ostrategy --budget 18                      # PASS
\end{verbatim}
}

\noindent Weak-game and strong $X$ certificates regenerated with \file{solver.cpp} may differ in
size from the included ones, which were exported by an earlier version of that solver; every
regenerated certificate passes the checkers. The export of \file{strong/solver\_strong.cpp} is
deterministic and reproduces the included $O$ strategies byte for byte. Apple clang~21 at plain
\texttt{-O2} miscompiles a range-based loop of this solver's placement generator (in the loop
vectorizer; the wrong placement set is caught by the solver's own symmetry self-check). The loop is
written with indices, and the solver and \file{strong/check6b.cpp} are built with
\texttt{-O2 -fno-vectorize} (with GCC, \texttt{-O2 -fno-tree-vectorize}).

\section*{Use of AI tools}

The programs in the ancillary files---all three solvers, all checkers and the driver scripts---and
drafts of this note were developed in sessions assisted by large-language-model coding tools.
The author checked the arguments and the proof-critical outputs and takes responsibility for the
note. The proofs do not depend on
the correctness of the search programs: they rest on the lemmas proved above and on certificates
that anyone can recheck with the included checkers. Agreement between separately written,
AI-assisted programs is an error filter, not an independent reproduction by other authors.

\begin{thebibliography}{99}
\raggedright

\bibitem{BCG82}
E.~R. Berlekamp, J.~H. Conway and R.~K. Guy,
\emph{Winning Ways for your Mathematical Plays}, Vol.~2, Academic Press, London, 1982,
pp.~684--685. (Not seen by the author; cited after~\cite{Tho96,Wiki}.)

\bibitem{EM22}
M.~Earnest and D.~Mathias, answer to ``Harary's generalized Tic-Tac-Toe; Winning strategy for
Skinny on a $7\times 7$ board?'', Puzzling Stack Exchange, January~15, 2022,
\url{https://puzzling.stackexchange.com/a/114453/3896}. (Not seen by the author; cited
after~\cite{Wiki}.)

\bibitem{Gar79}
M.~Gardner, Mathematical Games: In which players of ticktacktoe are taught to hunt bigger game,
\emph{Scientific American} \textbf{240}(4) (April 1979). (Not seen by the author; values quoted
after~\cite{Wiki}.)

\bibitem{Gar92}
M.~Gardner, \emph{Fractal Music, Hypercards and More\ldots: Mathematical Recreations from Scientific
American Magazine}, W.~H. Freeman, New York, 1992. (Table seen as reproduced at
\url{https://mathlair.allfunandgames.ca/tictactoe.php}.)

\bibitem{Gar01}
M.~Gardner, Harary's generalized ticktacktoe, Chapter~37 in \emph{The Colossal Book of
Mathematics}, W.~W. Norton, New York, 2001, pp.~493--503.

\bibitem{HSP07}
I.~Halupczok and J.-C. Schlage-Puchta, Achieving Snaky, \emph{Integers} \textbf{7} (2007), \#G02.

\bibitem{Nag02}
A.~Nagai, \emph{Df-pn Algorithm for Searching AND/OR Trees and Its Applications}, PhD thesis,
University of Tokyo, 2002.

\bibitem{OAI26}
OpenAI, Snaky in 21 Maker moves, preprint, September~25, 2026,
\url{https://github.com/openai/math/blob/main/preprints/Snaky-in-21-Maker-moves-September-25-2026/article.pdf}.

\bibitem{OEIS521}
OEIS Foundation Inc., The On-Line Encyclopedia of Integer Sequences, sequence A246521,
\url{https://oeis.org/A246521}, accessed October 7, 2026.

\bibitem{OEIS597}
P.~von Br\"omssen, Sequence A380597, in \emph{The On-Line Encyclopedia of Integer Sequences},
\url{https://oeis.org/A380597}, accessed October 7, 2026.

\bibitem{OEIS598}
P.~von Br\"omssen, Sequence A380598, in \emph{The On-Line Encyclopedia of Integer Sequences},
\url{https://oeis.org/A380598}, accessed October 7, 2026.

\bibitem{Sie04}
N.~Sieben, Snaky is a 41-dimensional winner, \emph{Integers} \textbf{4} (2004), \#G5.

\bibitem{Tho96}
C.~Thompson, Harary's animal-achievement game -- some questions, posting to \texttt{sci.math},
October 14, 1996, archived at \url{https://ics.uci.edu/~eppstein/junkyard/animal-game.html}.

\bibitem{Wiki}
Wikipedia, Harary's generalized tic-tac-toe,
\url{https://en.wikipedia.org/wiki/Harary's_generalized_tic-tac-toe}, accessed October 7, 2026.

\end{thebibliography}

\end{document}
